\section{Evaluation Protocol}
\label{sec:protocol}

A benchmark is only as trustworthy as its evaluation protocol. \wcabench{}
rests on four pillars: leakage prevention, stratified reporting, statistical
rigour, and reproducibility. Empirical rigour of this kind is increasingly
recognized as essential for progress \citep{sculley_rigor}. This section
specifies each pillar in turn.

\subsection{Leakage prevention}
\label{sec:protocol-leakage}

\wcabench{} uses a rolling-window evaluation that simulates deployment: for
every test competition, only data that existed \emph{before} that competition
may influence a prediction. Formally, for a target competition with date $d$,
the admissible information set is
$\mathcal{F}_d = \{\, r : \mathrm{date}(r) < d \,\}$, and no record at or after
$d$---including other competitors' results in the same round---may be used.

Three mechanisms enforce this. First, every feature function must accept an
explicit \asof{} timestamp; a feature that cannot be evaluated ``as of'' a
moment is not admissible. Feature caches are keyed by \asof{} so results cannot
be silently reused across time. Second, a testing utility
\code{assert\_no\_leakage(features, target\_date)} asserts that no feature
depends on a record dated at or after the target. Third, benchmark statistics
are frozen, as described next.

\subsection{Frozen statistics}
\label{sec:protocol-frozen}

Statistics that would otherwise drift as the test window advances are computed
once from the training window and held fixed for the entire test period.
Table~\ref{tab:frozen} lists them. Freezing is essential: if, for example,
``current world record'' were recomputed per test competition, a model could
benefit from information that a real deployment would not have at the decision
moment. The export contains \textbf{612{,}224} frozen person--event
statistics, each computed on the training split only.

\begin{table}[t]
  \centering
  \small
  \begin{threeparttable}
  \caption{Statistics frozen from the training window and reused for the whole
  test window.}
  \label{tab:frozen}
  \begin{tabular}{@{}lll@{}}
    \toprule
    Statistic & Computed from & Frozen for \\
    \midrule
    Competitor historical mean / variance & training ($\le 2022$) & whole test window \\
    Event world record & training & whole test window \\
    Event reference distribution & training & whole test window \\
    Skill-level percentile thresholds & training & whole test window \\
    \bottomrule
  \end{tabular}
  \end{threeparttable}
\end{table}

\subsection{Rolling-window algorithm}
\label{sec:protocol-algorithm}

Algorithm~\ref{alg:rolling} gives the protocol in pseudocode. Test competitions
are processed in chronological order; features are built from the frozen
statistics and the running history, a leakage assertion is checked, and
per-competition reports are aggregated at the end.

\begin{algorithm}[t]
\caption{Rolling-window evaluation with frozen statistics.}
\label{alg:rolling}
\begin{algorithmic}[1]
\Require model $f$, test competitions $C$, frozen statistics $F$
\Ensure aggregated report
\State $R \gets [\,]$
\State $C \gets \mathrm{sort}(C, \text{key} = \text{date})$ \Comment{chronological}
\ForAll{$c \in C$}
    \State $\x \gets \textsc{Featurize}(\asof = c.\text{date}, F)$
    \State \textbf{assert} \Call{NoLeakage}{$\x$, $c.\text{date}$}
    \State $\hat{\y} \gets f.\mathrm{predict}(\x)$
    \State $R.\mathrm{append}\big(\textsc{Evaluate}(\hat{\y}, c.\text{truth})\big)$
\EndFor
\State \Return \Call{Aggregate}{$R$}
\end{algorithmic}
\end{algorithm}

\subsection{Four-dimensional stratification}
\label{sec:protocol-stratification}

Aggregate metrics can be dominated by a single large subgroup. To prevent
conclusions from being masked, \wcabench{} mandates stratified reports along
four dimensions, summarized in Table~\ref{tab:strat}. All four are emitted
together with confidence intervals, so a reader can see whether a headline gain
holds within important subgroups.

\begin{table}[t]
  \centering
  \small
  \begin{threeparttable}
  \caption{The four stratification dimensions required for every task.}
  \label{tab:strat}
  \begin{tabular}{@{}lll@{}}
    \toprule
    Dimension & Groups & Rationale \\
    \midrule
    Event & speed-solving / variant / blind / other & data volume differs by orders of magnitude \\
    Skill level & novice (0--25\%), intermediate (25--75\%), & behaviour and stability differ by level \\
                & advanced (75--95\%), elite ($>$95\%) & \\
    Time slice & Test-A / Test-B / Test-C & temporal robustness \\
    Region & Asia, Europe, North America, & cross-regional generalization \\
           & South America, Oceania, Africa & \\
    \bottomrule
  \end{tabular}
  \end{threeparttable}
\end{table}

\paragraph{Hard subsets.}
To guard against benchmark saturation, each task also reports a hard subset:
for \task{3}, competitors whose historical \DNF{} rate is in $[0.1, 0.3]$; for
\task{1}, rounds in the top quartile of recent-result variance; for \task{2},
rounds in which the top-8 competitors differ by less than a threshold; and for
\task{5}, event pairs in the bottom quartile by sample size.

\paragraph{Cold start.}
Competitors whose first appearance falls inside the test window are flagged and
reported separately; they never enter the headline metric.

\subsection{Statistical testing}
\label{sec:protocol-stats}

Every claim of improvement must survive statistical scrutiny. \wcabench{}
prescribes the following tests, all reported with effect sizes.

\begin{table}[t]
  \centering
  \footnotesize
  \begin{threeparttable}
  \caption{Required statistical tools. Paired units are
  (competition $\times$ event $\times$ round).}
  \label{tab:stats}
  \begin{tabularx}{\linewidth}{@{}>{\raggedright\arraybackslash}p{3.3cm}
                                >{\raggedright\arraybackslash}X
                                >{\raggedright\arraybackslash}X@{}}
    \toprule
    Method & Purpose & Applicability \\
    \midrule
    Paired $t$-test & compare two models on paired errors
      & two models, paired samples \\
    Bootstrap CI ($1000\times$) & uncertainty for arbitrary metrics
      & any metric, unknown distribution \\
    Friedman $+$ Nemenyi & multi-model comparison, critical-difference diagram
      & $\ge 3$ models \\
    Effect size (Cohen's $d$ / Cliff's $\delta$)
      & distinguish significance from importance & all comparisons \\
    \bottomrule
  \end{tabularx}
  \end{threeparttable}
\end{table}

\paragraph{Paired $t$-test.}
We compare two models on the same test instances by testing the paired
difference of their per-instance errors, using a single pairing unit
(competition $\times$ event $\times$ round). We report the mean difference, a
95\% confidence interval, the $p$-value and an effect size.

\paragraph{Bootstrap confidence intervals.}
Algorithm~\ref{alg:bootstrap} gives the bootstrap used for arbitrary metrics.
The random seed is fixed (default 42) and recorded, and the number of
resamples (default 1000) is reported so that confidence intervals are exactly
reproducible.

\begin{algorithm}[t]
\caption{Bootstrap confidence interval for a metric.}
\label{alg:bootstrap}
\begin{algorithmic}[1]
\Require values $V$, metric $M$, resamples $B = 1000$, level $\alpha = 0.05$, seed $s = 42$
\Ensure point estimate and CI
\State $\mathrm{rng} \gets \textsc{DefaultRng}(s)$
\State $S \gets [\,]$
\For{$b = 1$ \textbf{to} $B$}
    \State $V_b \gets \textsc{SampleWithReplacement}(V, |V|, \mathrm{rng})$
    \State $S.\mathrm{append}(M(V_b))$
\EndFor
\State $(\text{lo}, \text{hi}) \gets \mathrm{percentile}(S,\ [50\alpha,\ 50(2-\alpha)])$
\State \Return $(M(V), \text{lo}, \text{hi})$
\end{algorithmic}
\end{algorithm}

\paragraph{Multi-model comparison.}
When three or more models are compared, a Friedman test first establishes
whether a difference exists; if it does, a Nemenyi post-hoc test performs
pairwise comparisons and a critical-difference diagram summarizes the result.
Multiple comparisons are corrected (Bonferroni or Nemenyi) to control
false-positive inflation.

\paragraph{Effect sizes.}
Significance does not imply importance, so we require Cohen's $d$ (0.2 small /
0.5 medium / 0.8 large) or Cliff's $\delta$ ($|\delta|<0.147$ negligible,
$<0.33$ small, $<0.474$ medium, otherwise large) for every comparison.

\paragraph{Common pitfalls.}
The protocol explicitly forbids comparing unpaired aggregate means, omitting
effect sizes, leaving bootstrap seeds unrecorded, and ignoring stratification.
Each of these would inflate apparent improvements.

\subsection{Reproducibility}
\label{sec:protocol-repro}

Reproducibility is graded and enforced. Submissions must provide training code
and seeds, the preprocessing script or a pipeline version reference, model
weights, and a compute-cost report. All random sources (Python, NumPy, deep
learning frameworks) are fixed, and headline results report the mean and
standard deviation over at least three seeds. We define three reproducibility
levels: \textbf{L1} (code reruns), \textbf{L2} (L1 plus results within a
$\pm 1\%$ metric tolerance) and \textbf{L3} (L2 plus downloadable raw
predictions and independently recomputable metrics). The minimum bar for
inclusion is L2; the main leaderboard requires L3. Preprocessing artifacts
record row counts and SHA-256 checksums for verification.

\paragraph{Report schema.}
Every model submission emits a structured report containing overall metrics,
the four stratified views, calibration and bootstrap intervals, significance
tests against baselines, and a compute-cost file. This uniformity is what makes
``the same evaluation code scores every model'' a reality rather than a slogan.
